\chapter*{Editor's note}

{\itshape\setlength{\parskip}{1.2ex}\setlength{\parindent}{0pt}%
These are notes for Module 1 of the ``Introduction to String Theory’’ course at the University of Bologna. Apart from this page,
they were written entirely by Claude Opus 5.5 from scans of handwritten lecture notes taken by me and other students in 2026.
They are an independent reconstruction, not a record of the lectures.
The lecturer has neither seen nor endorsed them, and they may depart from the lectures in content, emphasis and conventions.
Any errors are mine (and Claude’s).

I hope these notes prove useful, although I doubt they will hold much value in the years to come, when far more powerful
models will be available.

This leads me to some (somewhat uncalled-for and unprofessional) reflections on the strange times
I am living through. I remember the moment, a few years ago, when a new OpenAI model turned out to be a better coder than me.
Given how fast these models were improving, I knew I would never catch up.
Now I feel the same way about physics. I find it hard to believe that I will ever solve a physics problem that an AI
could not solve faster and more cheaply. This has made me genuinely sad for months now.

But, to complicate matters, mixed in with that sadness is real excitement. What will these systems teach us
about the deepest mysteries in physics? About quantum gravity and string theory? Which problems
will they solve? Perhaps the most important practical question a young graduate student can now ask is: ``How can I
position myself to understand the output of the equivalent of a thousand Wittens, working nonstop at fifty times human speed?''
I hope these lecture notes can help with that task.
 \par}

\par\bigskip
\begin{flushright}
  \textit{LazyLagrangian}\\
  \textit{6 November 2026}
\end{flushright}
